\section{Недоопределённость опыта и условный мост}
\label{sec:boundary}
Сохранение бита устанавливает определённое свойство системы при заданных вмешательствах. Можно ли из этой совокупности данных вывести, что данный экземпляр в данный момент испытывает субъективное переживание?

\paragraph{T5.} Пусть все аксиомы операциональной теории $T$ формулируются без предиката опыта $E(a,t)$. Возьмём любую модель $M$ этой теории с выделенной точкой $u=(a,t)$. Расширим её один раз пустой интерпретацией $E$, во второй раз --- интерпретацией, содержащей $u$. Поскольку меняется только новый предикат, обе структуры удовлетворяют $T$ и совпадают на исходах всех операциональных вмешательств. В первой $E(u)$ ложно, во второй истинно. Таким образом, $T$ не выводит ни одно из двух значений даже при выполненном функциональном критерии пересмотра. Это свойство выразимости теории. Построение двух структур не устанавливает физическую реализуемость двух различных переживаний на одном носителе.

\paragraph{T6.} Пакет $B=(B_{\mathrm{id}},B_{\mathrm{nom}},B_{\mathrm{epist}})$ имеет статус \texttt{DECLARED\_NOT\_DERIVED}. Его три части определяют носитель вывода, дополнительную посылку и правила обращения со свидетельством.
\begin{description}[leftmargin=1.35cm,style=nextline]
\item[$B_{\mathrm{id}}$] Задаются граница системы, макросостояние и допустимая группа переименований координат. Описание сохраняется при этих переименованиях. При фиксированном внешнем входе состояние должно влиять на преемника по внутренней причинной цепи. Для непрерывной линии требуется единственность преемника. При симметричном клонировании KV и независимом продолжении обеих копий регистрируется \texttt{BranchEvent}. У двух причинных потомков общий предок и разные идентификаторы ветвей. Это соглашение об учёте историй. Численное тождество субъекта остаётся открытым.
\item[$B_{\mathrm{nom}}$] Обозначим $C=I_{\mathrm{cont}}\land C_{\mathrm{med}}\land A_{\mathrm{ind}}\land D_{\mathrm{dyn}}$. Здесь $I_{\mathrm{cont}}$ задаёт непрерывность, $C_{\mathrm{med}}$ --- внутреннее посредничество при закрытых позициях прошлых выбора и сигнала с переносом того же правила на иной тип задачи. Для $A_{\mathrm{ind}}$ требуются независимые проверки доступа нескольких модулей к общему содержанию и изменения действия по оценке надёжности. Для $D_{\mathrm{dyn}}$ нужна устойчивая реорганизация, вызванная обработкой и причинно меняющая дальнейшую обработку после эпизода. Принимаем в качестве гипотезы достаточное условие $\forall u\,[C(u)\Rightarrow E(u)]$.
\item[$B_{\mathrm{epist}}$] Статус \texttt{UN} обозначает открытость вопроса при имеющихся свидетельствах. Управленческие метки \texttt{NONE}, \texttt{INVESTIGATION\_PRIORITY} и \texttt{SENTIENCE\_CANDIDATE} требуют отдельно указанных оснований. Для численной степени доверия нужны априорные вероятности и модель наблюдений.
\end{description}
\begin{theorem}[Условная теорема T6]
При принятии $B_{\mathrm{nom}}$ из $C(u)$ следует $E(u)$. Если существует хотя бы одна операциональная модель $M$, совместимая со свидетельствами и $B_{\mathrm{id}}$, в которой $C(u)$ ложно, то те же свидетельства после добавления $B_{\mathrm{nom}}$ совместимы с обоими значениями $E(u)$.
\end{theorem}
\begin{proof}
Первое утверждение есть применение импликации в точке $u$. Для второго положим $S=\{v:C(v)\}$ в $M$ и расширим одну и ту же операциональную модель интерпретациями $E_0=S$ и $E_1=S\cup\{u\}$. Обе включают $S$ и потому выполняют мостовой постулат. Поскольку $u\notin S$, они расходятся в значении $E(u)$. Операциональные данные и результаты всех заданных вмешательств у расширений общие.
\end{proof}
Отсутствие измерения конъюнкта само по себе не доказывает его ложность. Для Rule 110, однобитного регистра, табличного декодера и замороженной OLMo операциональный статус $E(a,t)$ в этой работе сохраняется как \texttt{UN}. Перенос внешнего токена оставляет открытыми внутреннее посредничество и независимые архитектурные индикаторы. Изменение KV внутри одного эпизода также недостаточно для вывода об устойчивой пластичности.
